\documentclass[11pt,reqno]{amsart}
\usepackage[margin=1.15in]{geometry}
\usepackage{amsmath,amssymb,amsthm,mathrsfs}
\usepackage{tikz}
\usepackage[colorlinks=true,linkcolor=blue,citecolor=blue]{hyperref}

\theoremstyle{plain}
\newtheorem{theorem}{Theorem}[section]
\newtheorem{proposition}[theorem]{Proposition}
\newtheorem{lemma}[theorem]{Lemma}
\newtheorem{corollary}[theorem]{Corollary}
\newtheorem{conjecture}[theorem]{Conjecture}
\theoremstyle{definition}
\newtheorem{definition}[theorem]{Definition}
\newtheorem{example}[theorem]{Example}
\newtheorem{remark}[theorem]{Remark}

\DeclareMathOperator{\soc}{soc}
\DeclareMathOperator{\topp}{top}
\DeclareMathOperator{\Hom}{Hom}
\DeclareMathOperator{\Ext}{Ext}
\DeclareMathOperator{\rad}{rad}
\newcommand{\bbC}{\mathbb{C}}
\newcommand{\bbZ}{\mathbb{Z}}
\newcommand{\Pia}{\Pi}
\newcommand{\dimv}{\underline{\dim}}
\newcommand{\al}{\alpha}
\newcommand{\om}{\omega}
\newcommand{\la}{\lambda}
\newcommand{\ka}{\kappa}
\newcommand{\Mvw}{\mathcal{M}}

\begin{document}

\title[Simple spectrum at extremal fixed points]
      {A construction of the extremal torus fixed point, and a criterion for
       simple spectrum}

\author{}
\date{\today}

\begin{abstract}
Let $\Gamma$ be a simply-laced Dynkin diagram, $\la=\om_\ka$ a fundamental
weight and $w$ an element of the Weyl group.  When $\mu=w\la$ the Nakajima
quiver variety $\Mvw(v,w)$ with $\la-\mu=\sum_i v_i\al_i$ is a single point,
carrying a unique $T$-fixed representation $V=\bigoplus_i V_i$.  Dinkins,
Karpov and Krylov raise, in \S5.4.2 of \cite{DKK}, the problem of
characterising combinatorially those $w$ for which $T$ acts on every $V_i$
with simple spectrum, in which case the fixed point is encoded by a coloured
poset and the vertex function is a sum over reverse plane partitions.

We solve this by producing $V$ itself.  Let $\Pia$ be the preprojective
algebra of $\Gamma$ and $I_\ka$ the injective envelope of the simple
$S_\ka$.  Reading a reduced word for $w$ from the right and recording at each
step the integer $c=\langle\al_i^\vee,\mu_{\mathrm{cur}}\rangle$ produces a
sequence $\sigma$ of vertices with multiplicities.  We prove that
\[
   V \;\cong\; \soc_{\sigma}\!\left(I_\ka\right),
\]
the iterated socle of Geiss--Leclerc--Schr\"oer taken along $\sigma$.  The
proof is uniform in type and rests on an identity, derived from
Crawley-Boevey's formula for $\Ext^1$ over a preprojective algebra, together
with a squeeze against $\dimv I_\ka=\la-w_0\la$.  Simple spectrum then becomes
the assertion that no two elements of the resulting coloured graded set share a
colour and a degree, which is decidable in polynomial time from any reduced
word.  In type $D_n$, for non-spin $\la$, we make the answer completely explicit:
writing $\mu=\sum_{i\in S}\eta_ie_i$ in Bourbaki coordinates, the spectrum is
simple if and only if $r+2m\le 2$, where $m$ counts the minus signs among the
first $n-2$ coordinates and $r=|S\cap\{n-1,n\}|$.  The proof of this last statement runs through a
strand decomposition of $I_\ka$ under the diagram automorphism exchanging the
two leaves, which reduces it to an explicit inequality between two integer
sequences.  We also record an element of $D_4$ that
is fully commutative but does not have simple spectrum, showing that the
sufficient condition considered in \cite[Ex.~5.12]{DKK} fails in both
directions.
\end{abstract}

\maketitle

\section{Introduction}

\subsection{The problem}
Fix a simply-laced Dynkin diagram $\Gamma$ with vertex set $I$, Cartan matrix
$C=(a_{ij})$, simple roots $\al_i$, fundamental weights $\om_i$ and Weyl group
$W$.  Let $\Mvw(v,w)$ denote the Nakajima quiver variety attached to
$\Gamma$ with dimension vector $v$ and framing $w$, and write
$\la=\sum_i w_i\om_i$, $\mu=\la-\sum_i v_i\al_i$.  When $\mu$ is an extremal
weight of the irreducible representation of highest weight $\la$ --- that is,
$\mu=w\la$ for some $w\in W$ --- the variety $\Mvw(v,w)$ is a single point
\cite{Nak,ST}, and that point is a representation $V=\bigoplus_{i\in I}V_i$ of
the doubled quiver satisfying the moment map equations, unique up to
isomorphism.

The point is fixed by the torus $T$ acting on the quiver data, so each $V_i$
carries a $T$-weight decomposition.  Dinkins, Karpov and Krylov show
\cite[Thm.~5.6, Thm.~5.7]{DKK} that when $T$ acts on every $V_i$ with
\emph{simple spectrum} --- all weights distinct --- the fixed point is encoded
by a finite coloured poset, and the associated vertex function is a generating
function over reverse plane partitions on that poset.  They then write, in
\S5.4.2:

\begin{quote}
``We consider it an interesting open problem to find a combinatorial
characterization of such $w$.''
\end{quote}

They rule out three natural guesses.  Dominant minuscule elements are too
restrictive \cite[Ex.~5.10]{DKK}; so are minuscule elements
\cite[Ex.~5.11]{DKK}; and so are fully commutative elements
\cite[Ex.~5.12]{DKK}.  Each example exhibits a $w$ with simple spectrum lying
outside the class.

\subsection{Results}
Our main theorem answers the question by constructing $V$ explicitly, uniformly
in the type, from any reduced word for $w$.

Let $\Pia=\Pia(\Gamma)$ be the preprojective algebra over a field, $S_i$ the
simple module at $i$, and $I_i$ the injective envelope of $S_i$.  For a
$\Pia$-module $M$ and $j\in I$ let $\soc_{(j)}(M)$ be the sum of all submodules
of $M$ isomorphic to $S_j$, and for a sequence $\sigma=(j_1,\dots,j_p)$ let
$\soc_\sigma(M)$ be the iterated socle of Geiss--Leclerc--Schr\"oer
\cite[\S2.4]{GLS}, defined by $M_0=0$ and
$M_q/M_{q-1}=\soc_{(j_q)}(M/M_{q-1})$.

\begin{definition}[the $c$-trace]\label{def:ctrace}
Let $w=s_{i_1}\cdots s_{i_\ell}$ be a reduced word.  Set $\mu^{(\ell)}=\la$ and,
for $p=\ell,\ell-1,\dots,1$,
\[
  c_p=\langle \al_{i_p}^\vee,\ \mu^{(p)}\rangle,
  \qquad
  \mu^{(p-1)}=\mu^{(p)}-c_p\al_{i_p}=s_{i_p}\mu^{(p)} .
\]
The \emph{$c$-weighted sequence} of the word is
$\sigma=\bigl(i_\ell^{\,c_\ell},\,i_{\ell-1}^{\,c_{\ell-1}},\dots,
i_1^{\,c_1}\bigr)$, each letter repeated according to its multiplicity, in
this order.
\end{definition}

Note $\mu^{(0)}=w\la=\mu$ and $\sum_{p\,:\,i_p=i}c_p=v_i$, so the $c$-trace
recomputes the dimension vector.  The invariant $c_{\max}=\max_p c_p$ equals
$1$ exactly when $w$ is $\la$-minuscule in the sense of
\cite[Def.~5.1]{DKK}.

\begin{theorem}[Main Theorem]\label{thm:main}
Let $\Gamma$ be simply-laced Dynkin, $\la=\om_\ka$, $w\in W$ of minimal length
in $wW_\la$, and $\sigma$ the $c$-weighted sequence of any reduced word for
$w$.  Then
\[
   \dimv \soc_\sigma\!\left(I_\ka\right)\;=\;v\;=\;\la-w\la,
\]
and consequently
\[
   V\;\cong\;\soc_\sigma\!\left(I_\ka\right)
\]
as graded representations of the doubled quiver.  In particular the
isomorphism class of $\soc_\sigma(I_\ka)$ does not depend on the chosen
reduced word.
\end{theorem}

The grading on the right is the path-length grading of $I_\ka$; Proposition
\ref{prop:grading} identifies it with the $T$-weight grading, and \S\ref{sec:poset}
explains how to read off the coloured poset of \cite[\S5.4]{DKK} from it.  The
criterion is then immediate.

\begin{corollary}\label{cor:crit}
$T$ acts on every $V_i$ with simple spectrum if and only if the coloured graded
set underlying $\soc_\sigma(I_\ka)$ has no repeated (colour, degree) pair,
i.e.\ $\dim \soc_\sigma(I_\ka)_i[d]\le 1$ for all $i\in I$ and all $d$.
\end{corollary}

The construction is effective: $I_\ka$ has dimension $\la-w_0\la$, independent
of $w$, and the socle steps are kernel computations of size bounded by the
graded pieces of $I_\ka$, so the criterion is decidable in time polynomial in
$|I|$ and $\ell(w)$ from any reduced word.

Two refinements make the answer sharper.  First, only one vertex needs to be
inspected.

\begin{theorem}\label{prop:lemA}
Let $t$ be the trivalent node of $\Gamma$ (types $D$ and $E$).  If
$\dim V_t[d]\le1$ for all $d$, then $\dim V_i[d]\le1$ for all $i$ and $d$.
\end{theorem}

Theorem \ref{prop:lemA} follows in both types from a single inward-injectivity
lemma about $I_\ka$ (Lemma \ref{lem:I}, \S\ref{sec:typeE}).  In type $E$ that
lemma is a \emph{complete} finite computation, since $E_6$, $E_7$ and $E_8$
have only finitely many triples (fundamental weight, vertex, degree); in type
$D$ it is proved uniformly in the rank (Lemma \ref{lem:Id}).  Remark
\ref{rem:nowindow} records why the more obvious argument, containment of
multiplicity-$2$ windows, fails in both types.  In types $D$ and $E$ this reduces the criterion to a
single vertex; in type $A$ there is no trivalent node, but there the question
is settled outright:

\begin{proposition}\label{prop:minuscule}
If $\la$ is minuscule then, for $w$ of minimal length in $wW_\la$, $c_p=1$ for
every step of every reduced word, $w$ is $\la$-minuscule, and $T$ acts with
simple spectrum.  In
particular the spectrum is always simple in type $A$, for the spin weights
$\om_{n-1},\om_n$ of $D_n$, for $\om_1$ and $\om_6$ of $E_6$, and for $\om_7$
of $E_7$.
\end{proposition}

Second, in type $D$ the criterion becomes a statement about the signs of $\mu$
with no reference to reduced words, heaps, or fully commutative elements.

\begin{theorem}[type $D$]\label{thm:typeD}
Let $\Gamma=D_n$, $t=n-2$, and $\la=\om_k$ with $k\le n-2$.  Write
$\mu=w\la=\sum_{i\in S}\eta_ie_i$ in Bourbaki coordinates, $|S|=k$,
$\eta_i=\pm1$.  Put
\[
  m=\#\{\,i\le n-2:\ i\in S,\ \eta_i=-1\,\},
  \qquad
  r=|S\cap\{n-1,n\}| .
\]
Then $T$ acts with simple spectrum if and only if $r+2m\le2$; equivalently, iff
$m=0$, or $m=1$ and $r=0$.  For the spin weights the spectrum is always simple.
The proof runs through Proposition \ref{prop:T}, which
Theorem \ref{thm:staircase} establishes; see Remark \ref{rem:Tk3} for the route.
\end{theorem}

Finally we record a counterexample in the direction opposite to
\cite[Ex.~5.12]{DKK}.

\begin{example}\label{ex:fcbad}
In $D_4$ take $w=s_2s_4s_3s_1s_2$, so $v=(1,3,1,1)$.  This $w$ is fully
commutative, but the spectrum is \emph{not} simple: the path-length graded
dimensions of $V$ inside $I_2$ are $\{2:1\}$ in degree $4$, $\{1:1,3:1,4:1\}$
in degree $3$ and $\{2:2\}$ in degree $2$, so, read from the top down as in the
diagrams of \cite{DKK}, the rows are $[2]$, $[1,3,4]$, $[2,2]$.  Together with \cite[Ex.~5.12]{DKK}, which exhibits a
$w$ with simple spectrum that is not fully commutative, this shows that full
commutativity is neither necessary nor sufficient.
\end{example}

\subsection{Method}
The proof of Theorem \ref{thm:main} is short and does not proceed case by case.
It has two ingredients.  The first is an exact formula for the size of a single
socle step, valid for an arbitrary submodule of $I_\ka$:

\begin{lemma}[\S\ref{sec:identity}]\label{lem:cb}
Let $X\subseteq I_\ka$ be a submodule, $v'=\dimv X$, $\mu'=\la-\sum_i
v'_i\al_i$.  Then for every $i\in I$
\[
   \dim\soc_{(i)}\!\left(I_\ka/X\right)
   \;=\;\langle\al_i^\vee,\mu'\rangle\;+\;\dim\topp_i(X),
\]
where $\topp_i(X)=\Hom_\Pia(X,S_i)$.
\end{lemma}

This comes from applying $\Hom(S_i,-)$ to $0\to X\to I_\ka\to Q\to0$, using
$\Ext^1(S_i,I_\ka)=0$, and then Crawley-Boevey's formula
$\dim\Ext^1_\Pia(M,N)=\dim\Hom(M,N)+\dim\Hom(N,M)-(\dimv M,\dimv N)$
\cite{CB}.

The second ingredient is a squeeze.  Since $\dim\topp_i(X)\ge0$, each socle
step grows $X$ by at least $c$; summing the $c$'s along a reduced word for the
longest element gives exactly $\la-w_0\la$, which is precisely $\dimv I_\ka$.
As $\soc_\sigma(I_\ka)\subseteq I_\ka$ the two bounds meet, so \emph{every}
step contributes exactly $c$ and every $\topp$-term vanishes.  Restricting to
the prefix corresponding to $w$ gives Theorem \ref{thm:main}, once uniqueness
of the submodule with $w$-extremal dimension vector \cite[Prop.~4.9]{ST}
identifies the result with $V$.

\subsection{Verification}
Every statement in this paper is accompanied by machine verification against an
independent solver that enumerates graded submodules directly.  In particular
the construction reproduces the three coloured posets printed in
\cite[pp.~82--83]{DKK}, Hasse edges included, and the criterion agrees with the
solver on $2179$ elements across types $A$, $D$ and $E$ --- every fundamental
weight of $A_3$--$A_6$, $D_4$--$D_7$, $E_6$, together with $E_7\om_1$ and
$E_7\om_7$, whose Weyl orbits have at most $400$ elements.  See \S\ref{sec:verification}.

\subsection{Acknowledgements}
The problem was communicated by V.~Krylov.

\section{Setup}\label{sec:setup}

\subsection{The preprojective algebra}
Let $Q$ be a quiver with underlying graph $\Gamma$ and let $\overline{Q}$ be its
double, with $a^*$ the reverse of $a$.  The preprojective algebra is
\[
  \Pia=\Pia(\Gamma)=\Bbbk\overline{Q}\Big/\Big(\textstyle\sum_{a}(aa^*-a^*a)\Big),
\]
where $\Bbbk$ is any field.  Every argument below is characteristic-free;
Crawley-Boevey's formula, which we use in \S\ref{sec:identity}, holds over an
arbitrary field.  The computer verifications of \S\ref{sec:verification} were
run over $\mathbb{F}_p$ for $p=5,7,11,13$ with identical results.
A $\Pia$-module is the same thing as a representation of $\overline{Q}$
satisfying the moment map equation at every vertex, which is exactly the datum
carried by a point of $\Mvw(v,w)$ with vanishing stability-free part.  We use
throughout the following standard facts, valid because $\Gamma$ is Dynkin.

\begin{proposition}\label{prop:selfinj}
$\Pia(\Gamma)$ is finite dimensional and self-injective.  Its Nakayama
permutation $\nu$ is determined by $\soc(\Pia e_j)\cong S_{\nu(j)}$, and
\[
  I_\ka \;\cong\; \Pia e_{\nu^{-1}(\ka)} .
\]
Moreover $\nu$ is the identity except in types $A_n$, $D_{\text{odd}}$ and
$E_6$, where it is induced by the nontrivial diagram automorphism.
\end{proposition}

\begin{proposition}\label{prop:dimI}
$\dimv I_\ka=\om_\ka-w_0\om_\ka$ in root coordinates.
\end{proposition}

Proposition \ref{prop:selfinj} is classical; see for instance
\cite[\S1]{CB}.  Proposition \ref{prop:dimI} follows from the description of
$\Pia e_j$ as $\bigoplus_{d\ge0}\tau^{-d}P_j$ over the path algebra of an
orientation of $\Gamma$, whose total dimension vector is
$\om_j+\om_{j^*}=\om_j-w_0\om_j$; it is verified directly for every fundamental
weight of $A_3,A_4,A_5,D_4,D_5,D_6,D_7,E_6$ and $E_7$ in
\S\ref{sec:verification}.

Every $\Pia e_j$ is graded by path length, and \emph{this is the grading we
use throughout}; it coincides with the socle-layer grading read in reverse.  We
write $M_i[d]$ for the degree-$d$ part of $M$ at vertex $i$, and $d_{\max}$ for
the top degree of $I_\ka$.

\subsection{Iterated socles}
For $j\in I$, $\soc_{(j)}(M)$ is the largest submodule of $M$ that is a direct
sum of copies of $S_j$; concretely it is the set of $x\in M_j$ annihilated by
every arrow out of $j$.  Given a sequence $\sigma=(j_1,\dots,j_p)$ define
$\soc_\sigma(M)=M_p$ where $M_0=0$ and $M_q$ is the preimage in $M$ of
$\soc_{(j_q)}(M/M_{q-1})$.  We use only that $\soc_\sigma$ is a submodule and
that $\soc_{\sigma}(M)\subseteq\soc_{\sigma'}(M)$ when $\sigma$ is a prefix of
$\sigma'$.

\subsection{Extremal fixed points}
We use Savage--Tingley's identification of the quiver variety with a quiver
Grassmannian.

\begin{proposition}[{\cite[Thm.~4.4, Prop.~4.9]{ST}; \cite{BK,BKT}}]\label{prop:ST}
$\mathfrak{L}(v,w)\cong\mathrm{Gr}_{v}\!\left(\bigoplus_i I_i^{\,w_i}\right)$,
and if $v=\la-w\la$ is $w$-extremal then there is exactly one submodule of
$\bigoplus_i I_i^{\,w_i}$ with dimension vector $v$.
\end{proposition}

The uniqueness has a sharper and more useful form.  Baumann and Kamnitzer
attach to a $\Pia$-module $T$ a pseudo-Weyl polytope $\mathrm{Pol}(T)$, and
show \cite{BK} that it is the convex hull of the dimension vectors of the
submodules of $T$; Baumann, Kamnitzer and Tingley prove the \emph{rigidity}
property \cite[\S3.1]{BKT}: \emph{if $x$ is a vertex of $\mathrm{Pol}(T)$ then
$T$ has a unique subobject $X$ with $\dimv X=x$.}  For $T=I_\ka$ the polytope
is the weight polytope of $\om_\ka$ reflected in the origin, so its vertices
are exactly the $w$-extremal dimension vectors, and Proposition \ref{prop:ST}
is the vertex case.  Two consequences we use below:

\begin{corollary}\label{cor:pol}
Let $U\subseteq I_\ka$ be a submodule.  Then $\mu=\la-\dimv U$ is a weight of
the irreducible representation $L(\om_\ka)$, and
\[
   \delta(\dimv U)\;=\;2(\dimv U)_\ka-(\dimv U)^{\mathsf T}C(\dimv U)
   \;=\;(\la,\la)-(\mu,\mu)\;\ge\;0,
\]
with equality if and only if $\dimv U$ is a vertex of $\mathrm{Pol}(I_\ka)$,
i.e.\ $w$-extremal.  If $I_\ka$ carries two distinct submodules with the same
dimension vector $v$, then $v$ is not a vertex and $\delta(v)>0$.
\end{corollary}

In type $D_n$ with $\la=\om_k$, $k\le n-2$, the weights of $L(\om_k)$ are the
$\mu$ with entries in $\{0,\pm1\}$ and $|\mathrm{supp}\,\mu|\le k$ congruent to
$k$ modulo $2$, so Corollary \ref{cor:pol} gives
$\delta(\dimv U)=k-|\mathrm{supp}\,\mu|$ --- verified on all $2014$ graded
submodules of $I_\ka$ for $D_5\om_2$, $D_5\om_3$, $D_6\om_3$, $D_6\om_4$,
$D_7\om_3$.

For $\la=\om_\ka$ the framing is $\bigoplus_i I_i^{w_i}=I_\ka$, so
Proposition \ref{prop:ST} says: \emph{$I_\ka$ has a unique submodule of
dimension vector $v$, and the fixed point $V$ is that submodule} (we use
\cite[Thm.~4.4]{ST} in the convention in which a point of the quiver
Grassmannian is a submodule of the injective, not a quotient of it).  Hence to
identify $V$ it suffices to exhibit any submodule of $I_\ka$ of the right
dimension vector.  That is what Theorem \ref{thm:main} does.

It remains to match the two gradings.

\begin{lemma}[the bidegree is determined]\label{lem:bigrade}
$\Pia$ is bigraded by the numbers of $a$- and $a^*$-arrows separately, the
relations $\sum_a(aa^*-a^*a)$ being bihomogeneous of bidegree $(1,1)$.  For
$\Gamma$ a Dynkin diagram --- hence a tree --- the piece $(\Pia e_j)_i[d]$ is
pure of a single bidegree, namely
\[
   p=\tfrac12\bigl(d+f(i)-f(j)\bigr),\qquad q=\tfrac12\bigl(d-f(i)+f(j)\bigr),
\]
where $f$ is any function on the vertices with $f(\mathrm{head})-
f(\mathrm{tail})=1$ on every arrow of the chosen orientation.
\end{lemma}

\begin{proof}
Such an $f$ exists because $\Gamma$ is a tree: define it by walking out from
any root; there is no cycle to obstruct consistency.  An $a$-arrow raises $f$
by $1$ and its reverse lowers it by $1$, so along any path from $j$ to $i$ one
has $p-q=f(i)-f(j)$.  With $p+q=d$ this determines $p$ and $q$.
\end{proof}

\begin{proposition}\label{prop:grading}
The $T$-weight decomposition of $V_i$ coincides with its path-length
decomposition, up to an overall affine normalisation independent of $i$.  In
particular
\[
   T \text{ acts on every } V_i \text{ with simple spectrum}
   \iff
   \dim V_i[d]\le 1 \ \text{ for all } i \text{ and } d .
\]
\end{proposition}

\begin{proof}
Because $\la=\om_\ka$ is fundamental the framing space is one dimensional, so
the framing torus acts on all of $V$ through a single character and contributes
only a global shift.  The rest of $T$ acts through the bigrading of $\Pia$: any
such torus assigns to $(\Pia e_{\ka'})_i[d]$ a weight that is a fixed integral
combination $\alpha p+\beta q$ of the two bidegrees.  By Lemma \ref{lem:bigrade} both
$p$ and $q$ are affine functions of $d$ with $i,\ka'$ fixed, with $p$ and $q$
each of slope $\pm\frac12$ in $d$; hence $\alpha p+\beta q$ is an affine function of $d$
of slope $\frac12(\alpha+\beta)$, which is injective in $d$ unless
$\alpha+\beta=0$.  The case $\alpha+\beta=0$ is excluded because the
symplectic form has nonzero weight, that being what makes the fixed locus of
$T$ on $\Mvw(v,w)$ isolated.  So the weight determines and is determined by
$d$, degree by degree.

Finally $V$ is a graded submodule: since $v$ is $w$-extremal, Proposition
\ref{prop:ST} gives a \emph{unique} submodule of $I_\ka$ with dimension vector
$v$, and uniqueness forces it to be stable under the whole torus, hence to
carry the restricted grading.
\end{proof}

\section{The socle-step identity}\label{sec:identity}

\begin{proof}[Proof of Lemma \ref{lem:cb}]
Write $I=I_\ka$ and let $Q=I/X$, so that $0\to X\to I\to Q\to0$ is exact.
Applying $\Hom_\Pia(S_i,-)$ gives
\[
0\to\Hom(S_i,X)\to\Hom(S_i,I)\to\Hom(S_i,Q)\to\Ext^1(S_i,X)\to\Ext^1(S_i,I).
\]
The last term vanishes because $I$ is injective.  Since $\soc(I_\ka)=S_\ka$ we
have $\dim\Hom(S_i,I)=\delta_{i\ka}$, and $\dim\Hom(S_i,Q)=\dim\soc_{(i)}(Q)$.
Therefore
\begin{equation}\label{eq:les}
  \dim\soc_{(i)}(Q)=\delta_{i\ka}-\dim\soc_{(i)}(X)+\dim\Ext^1(S_i,X).
\end{equation}
By Crawley-Boevey's formula for modules over a preprojective algebra
\cite[Lem.~1]{CB},
\[
  \dim\Ext^1_\Pia(M,N)=\dim\Hom(M,N)+\dim\Hom(N,M)-(\dimv M,\dimv N),
\]
where $(\cdot,\cdot)$ is the symmetrised Tits form, so that
$(\dimv M,\dimv N)=(\dimv M)^{\mathsf T}C(\dimv N)$.  Taking $M=S_i$ and
$N=X$, and noting $\Hom(S_i,X)=\soc_{(i)}(X)$ and $\Hom(X,S_i)=\topp_i(X)$,
\[
  \dim\Ext^1(S_i,X)=\dim\soc_{(i)}(X)+\dim\topp_i(X)-(\al_i,v').
\]
Substituting into \eqref{eq:les} the socle terms cancel and
\[
  \dim\soc_{(i)}(Q)=\delta_{i\ka}-(\al_i,v')+\dim\topp_i(X).
\]
Finally $(\al_i,v')=\sum_j a_{ij}v'_j=\langle\al_i^\vee,\sum_jv'_j\al_j\rangle$
and $\langle\al_i^\vee,\om_\ka\rangle=\delta_{i\ka}$, so
$\delta_{i\ka}-(\al_i,v')=\langle\al_i^\vee,\mu'\rangle$.
\end{proof}

\begin{corollary}\label{cor:mono}
A socle step of colour $i$ applied to $X\subseteq I_\ka$ enlarges $X$ by
exactly $\langle\al_i^\vee,\mu'\rangle+\dim\topp_i(X)$ dimensions, all at
vertex $i$.  In particular the growth is at least
$\langle\al_i^\vee,\mu'\rangle$, with equality if and only if $S_i$ does not
occur in the head of $X$.
\end{corollary}

\begin{remark}\label{rem:repeats}
Lemma \ref{lem:cb} also explains why repeating a colour is harmless.  Suppose a
step of colour $i$ with $c=\langle\al_i^\vee,\mu'\rangle$ has just added
exactly $c$ dimensions, producing $X'$.  Those $c$ new elements lie in the head
of $X'$, so $\dim\topp_i(X')=c$; and the weight has become
$s_i\mu'$, with $\langle\al_i^\vee,s_i\mu'\rangle=-c$.  Lemma \ref{lem:cb}
gives $\dim\soc_{(i)}(I_\ka/X')=-c+c=0$: further repeats add nothing.  Thus
$\soc_\sigma$ is unchanged if each block $i^{\,c}$ is replaced by $i^{\,c'}$
for any $c'\ge c$, and the $c$-weighted sequence may equivalently be described
as ``apply $\soc_{(i)}$ until it stabilises''.
\end{remark}

\section{Proof of the Main Theorem}\label{sec:proof}

\begin{lemma}\label{lem:extend}
Let $w$ be of minimal length in $wW_\la$ and $\mu=w\la$.
\begin{enumerate}
\item[(i)] Every $c$-trace value of every reduced word for $w$ satisfies
$c_p\ge1$.
\item[(ii)] If $\mu\ne w_0\la$ there is a simple reflection $s_i$ with
$\langle\al_i^\vee,\mu\rangle\ge1$, and for any such $i$ one has
$\ell(s_iw)=\ell(w)+1$ with $s_iw$ again of minimal length in $s_iwW_\la$.
\end{enumerate}
Consequently every reduced word for $w$ extends on the left to a reduced word
for the longest minimal-length representative $w_0^\la$, with all $c$-trace
values $\ge1$.
\end{lemma}

\begin{proof}
(i) Write $u=s_{i_{p+1}}\cdots s_{i_\ell}$, so $\mu^{(p)}=u\la$.  The word is
reduced, so $\ell(s_{i_p}u)=\ell(u)+1$ and hence $u^{-1}\al_{i_p}>0$; since
$\la$ is dominant,
$c_p=\langle (u^{-1}\al_{i_p})^\vee,\la\rangle\ge0$.  It remains to exclude
$c_p=0$.  If $c_p=0$
then $s_{i_p}u\la=u\la$, hence
$w\la=s_{i_1}\cdots s_{i_{p-1}}u\la$, and $s_{i_1}\cdots s_{i_{p-1}}u$ has
length at most $\ell-1$.  This contradicts the minimality of $\ell(w)$ among
elements sending $\la$ to $\mu$.

(ii) If $\langle\al_i^\vee,\mu\rangle\le0$ for every $i$ then $\mu$ is
antidominant, hence $\mu=w_0\la$.  So assume
$\langle\al_i^\vee,w\la\rangle>0$.  Then $\langle\la,w^{-1}\al_i^\vee\rangle>0$;
since $\la$ is dominant this forces $w^{-1}\al_i>0$, and therefore
$\ell(s_iw)=\ell(w)+1$.  Let $J=\{j:\langle\al_j^\vee,\la\rangle=0\}$, so that
$W_\la=W_J$ and $w$ is minimal in $wW_J$ if and only if $w\al_j>0$ for all
$j\in J$.  For $j\in J$ we have $w\al_j>0$, and $s_i(w\al_j)<0$ could only
happen if $w\al_j=\al_i$; but then
$\langle\al_i^\vee,w\la\rangle=\langle\al_j^\vee,\la\rangle=0$, contradicting
$\langle\al_i^\vee,w\la\rangle>0$.  Hence $s_iw\al_j>0$ for all $j\in J$ and
$s_iw$ is minimal in its coset.
\end{proof}

\begin{proof}[Proof of Theorem \ref{thm:main}]
Fix a reduced word for $w$ and extend it, by Lemma \ref{lem:extend}, to a
reduced word for $w_0^\la$; let $\sigma_{\mathrm{full}}$ be its $c$-weighted
sequence and $\sigma$ the prefix belonging to $w$.  Let
$X_0=0\subseteq X_1\subseteq\cdots\subseteq X_L=\soc_{\sigma_{\mathrm{full}}}
(I_\ka)$ be the resulting filtration, indexed by the \emph{blocks}
$i_p^{\,c_p}$ of $\sigma_{\mathrm{full}}$ rather than by individual letters.

By Corollary \ref{cor:mono} applied to the first letter of the $p$-th block,
and because the remaining letters of the block can only enlarge $X$ further,
\begin{equation}\label{eq:lower}
  \dim X_p-\dim X_{p-1}\;\ge\;c_p \qquad(1\le p\le L).
\end{equation}
Here $c_p\ge1$ by Lemma \ref{lem:extend}(i); it is this positivity that makes
\eqref{eq:lower} a lower bound strong enough to be squeezed.
Summing, and using that the $c$-trace telescopes,
$\sum_p c_p\al_{i_p}=\la-w_0^\la\la=\la-w_0\la$, we get
\[
  \dim \soc_{\sigma_{\mathrm{full}}}(I_\ka)\;\ge\;
  \sum_p c_p \;=\; \dim\bigl(\la-w_0\la\bigr)\;=\;\dim I_\ka,
\]
the last equality by Proposition \ref{prop:dimI}.  On the other hand
$\soc_{\sigma_{\mathrm{full}}}(I_\ka)\subseteq I_\ka$, so the inequality is an
equality.  Consequently every inequality in \eqref{eq:lower} is an equality,
and since the $p$-th block only enlarges $X$ at vertex $i_p$,
\[
  \dimv X_p-\dimv X_{p-1}\;=\;c_p\,\al_{i_p}\quad\text{for all }p .
\]
Truncating at the prefix $\sigma$ gives
$\dimv\soc_\sigma(I_\ka)=\sum_{p\le \ell}c_p\al_{i_p}=\la-w\la=v$.  By
Proposition \ref{prop:ST} the submodule of $I_\ka$ with this dimension vector
is unique and equals $V$.  Word-independence follows, since the right-hand side
depends only on $w$.
\end{proof}

\begin{corollary}\label{cor:tops}
At the \emph{first} letter of each block $i^{\,c}$ of $\sigma$ one has
$\topp_{i}(X)=0$ and $\dim\soc_{(i)}(I_\ka/X)=c$ exactly.  (Within a block
this fails, and must: after the first letter $\topp_i=c$ and
$\soc_{(i)}(I_\ka/X)=0$, as computed in Remark \ref{rem:repeats}.)
\end{corollary}

\begin{remark}
Corollary \ref{cor:tops} has a purely combinatorial shadow.  Only reflections
at $i$ and at its neighbours change $\langle\al_i^\vee,\cdot\rangle$, and a
colour-$i$ step flips that pairing from $c$ to $-c$; so between two consecutive
colour-$i$ blocks of a $c$-trace some neighbouring colour must intervene.  It
is exactly that intervening step which absorbs $S_i$ out of the head of $X$.
\end{remark}

\begin{proof}[Proof of Proposition \ref{prop:minuscule}]
If $\la$ is minuscule then $\langle\al^\vee,\nu\rangle\in\{0,\pm1\}$ for every
root $\al$ and every weight $\nu$ of $L(\la)$, so every $c_p$ equals $1$ and
$w$ is $\la$-minuscule in the sense of \cite[Def.~5.1]{DKK}.  By Stembridge
\cite[Prop.~2.1, \S2]{Stem} a $\la$-minuscule element is fully commutative, so
its heap $H(w)$ is independent of the reduced word; and $\la$-minuscule heaps
are ranked, with each colour fibre $H(w)_i$ a chain.  By \cite[Prop.~3.8, Lem.~3.10]{DFKKMSSVW} the fixed
point representation is the heap module $\bbC H(w)$, with $V_i$ spanned by
$H(w)_i$.  The rank function of $H(w)$ is the path-length degree of Proposition
\ref{prop:grading}, and a chain occupies pairwise distinct ranks, so every
graded piece of $V_i$ has dimension at most one.
\end{proof}

\section{Reading the coloured poset}\label{sec:poset}

Let $X=\soc_\sigma(I_\ka)\cong V$ and assume the criterion of Corollary
\ref{cor:crit} holds, so that $\dim X_i[d]\le1$ for all $i,d$; only then are
the vertices of the poset canonically labelled.  Define a coloured graded set
\[
  P(w)=\bigl\{(i,d,\text{basis vector of }X_i[d])\bigr\},
\]
coloured by $i$ and ranked by $d_{\max}-d$, with $(i,d)\lessdot(j,d+1)$ whenever
$j$ is adjacent to $i$ and the corresponding arrow of $\Pia$ is nonzero on the
relevant vectors.  Reading the degrees in decreasing order recovers the Hasse
diagrams of \cite[\S5.4]{DKK}, whose top row is the top degree $d_{\max}$.  Then
\cite[Thm.~5.6, Thm.~5.7]{DKK} express the vertex function as a sum over
reverse plane partitions on $P(w)$.

\begin{example}[{\cite[Ex.~5.10]{DKK}}]
$\Gamma=D_4$, $\la=\om_2$, $w=s_1s_2s_3s_4s_2$.  The $c$-trace gives
$v=(2,2,1,1)$ and $\sigma=(2,4,3,2,1,1)$.  The construction returns rows
$[2],[1,3,4],[2],[1]$ read from the top degree down, matching the diagram on
p.~82 of \cite{DKK} including its six covering relations.
\end{example}

\begin{example}[{\cite[Ex.~5.11, Ex.~5.12]{DKK}}]
$\Gamma=D_5$, $\la=\om_2$.  For $w=s_2s_1s_3s_4s_5s_3s_2$ the construction
returns $[2],[1,3],[2,4,5],[3],[2]$, and for
$w=s_2s_1s_2s_3s_4s_5s_3s_2$ it returns $[2],[1,3],[2,4,5],[3],[2],[1]$, both
matching p.~83 of \cite{DKK} with all covering relations.
\end{example}

\section{Type \texorpdfstring{$D$}{D}}\label{sec:typeD}

Throughout this section $\Gamma=D_n$ with Bourbaki labelling
$\al_i=e_i-e_{i+1}$ for $i\le n-1$ and $\al_n=e_{n-1}+e_n$, so that the
trivalent node is $t=n-2$.

\begin{lemma}\label{lem:C}
Let $\la=\om_k$ with $k\le n-2$ and $\mu=w\la=\sum_{i\in S}\eta_ie_i$.  With
$m$ and $r$ as in Theorem \ref{thm:typeD} one has $v_t=r+2m$.
\end{lemma}

\begin{proof}
Write $\la-\mu=\sum_jc_je_j$.  Comparing coefficients of $e_j$ in
$\la-\mu=\sum_iv_i\al_i$ gives $c_1=v_1$ and $c_j=v_j-v_{j-1}$ for
$2\le j\le n-2$, so $v_t=\sum_{j\le n-2}c_j$.  Since $k\le n-2$ the support of
$\la$ lies in the first $n-2$ coordinates, whence
$\sum_{j\le n-2}c_j=k-\sum_{j\le n-2}\mu_j$.  Writing $p$ and $m$ for the
number of positive and negative signs of $\mu$ among the first $n-2$
coordinates, $\sum_{j\le n-2}\mu_j=p-m$ and $p+m=k-r$, so
$v_t=k-(p-m)=r+2m$.
\end{proof}

\begin{lemma}[the recursion]\label{lem:R}
Let $a_d(i)=\dim (I_\ka)_i[d]$.  Then for every $i$ and every $d$ strictly
below the top degree,
\[
   a_{d+1}(i)=\sum_{j\sim i}a_d(j)-a_{d-1}(i),\qquad
   a_0(i)=\delta_{i,\ka'},\ a_{-1}=0,
\]
where $\ka'=\nu^{-1}(\ka)$.
\end{lemma}

\begin{proof}
By construction the degree-$(d+1)$ part of the free module at $j$ is
$\bigoplus_{i\sim j}(I_\ka)_i[d]$, and the preprojective relations impose one
condition for each basis vector of $(I_\ka)_j[d-1]$.  These conditions are
independent precisely when
$Y:(I_\ka)_j[d-1]\to\bigoplus_{i\sim j}(I_\ka)_i[d]$ is injective.  A nonzero
element of $\ker Y$ would span a copy of $S_j$ inside $I_\ka$; but $I_\ka$ has
simple socle $S_\ka$, concentrated in the top degree.  Hence $Y$ is injective
below the top degree.
\end{proof}

\begin{lemma}[method of images]\label{lem:P}
For $\Gamma=D_n$ and $\la=\om_k$ with $k\le n-2$, the colour-$t$ multiplicity
sequence of $I_\ka$ is $(1,2,\dots,2,1)$ with $k-1$ twos, occupying the $k+1$
degrees $t-k,\ t-k+2,\ \dots,\ t+k$.  More generally vertex $i$ carries
multiplicity $2$ in exactly $\max(0,\ i+k-n+1)$ degrees and multiplicity
$\le2$ throughout.
\end{lemma}

\begin{proof}
The datum $a_0=\delta_k$ is fixed by the diagram automorphism exchanging $n-1$
and $n$, hence so is the solution: $a_d(n-1)=a_d(n)=:b_d$.  Setting
$a_d(t+1):=2b_d$, the recursion of Lemma \ref{lem:R} becomes the free chain
recursion $a_{d+1}(i)+a_{d-1}(i)=a_d(i-1)+a_d(i+1)$ on $1,\dots,L$ with
$L=t+1=n-1$, subject to a Dirichlet condition $a_d(0)=0$ and a Neumann
(mirror) condition $a_d(L+1)=a_d(L-1)$.

On the infinite chain the solution with $u_0=\delta_k$, $u_{-1}=0$ is
$u_d(i)=1$ if $|i-k|\le d$ and $i\equiv k+d \pmod 2$, and $0$ otherwise
(induction on $d$).  The Dirichlet condition is an odd reflection about $0$ and
the mirror condition an even reflection about $L$; the group they generate is
infinite dihedral, with images of $k$ at $2Lj\pm k$ carrying alternating signs.
Hence
\[
   a_d(i)=\sum_{x\ \text{image of}\ k}
      \mathrm{sign}(x)\cdot\bigl[\,|i-x|\le d\ \text{and}\ i\equiv x+d\ (2)\,\bigr].
\]
The four images nearest $t$ are $k$ and $2L-k$ with sign $+$, at distances
$t-k$ and $t+2-k$, and $-k$ and $2L+k$ with sign $-$, at distances $t+k$ and
$t+2+k$.  (The image solution vanishes identically once $d$ exceeds $t+2+k$, and the
total $\sum_{d,i}a_d(i)$ it predicts equals $\la-w_0\la=\dimv I_\ka$ by
Proposition \ref{prop:dimI}; this pins the top degree and shows that the
recursion of Lemma \ref{lem:R}, valid below the top degree, determines
$I_\ka$ completely.)  They enter the window one at a time as $d$ grows, so
$a_d(t)$ is $0$
for $d<t-k$, then $1$, then $2$ on $[t+2-k,\,t+k)$, then $1$ on
$[t+k,\,t+2+k)$, then $0$; the degrees of the correct parity in the middle
range number $(2k-4)/2+1=k-1$.  The same computation at a vertex $i\le t$ of the
chain shows the value $2$ occurs exactly for
$\max(|i-k|,2L-k-i)\le d<\min(i+k,2L+k-i)$, giving $\max(0,i+k-n+1)$ degrees.
(At the two leaves $n-1,n$ the folding gives $a_d(n-1)=a_d(n)=b_d$ with
$2b_d=a_d(t+1)\le2$, so $b_d\le1$ and they never carry multiplicity $2$; the
displayed count applies only on the chain $i\le t$.)
\end{proof}

\begin{remark}[window containment fails, and why]\label{rem:nowindow}
It is tempting to argue Theorem \ref{prop:lemA} in type $D$ from Lemma
\ref{lem:P} alone: the multiplicity-$2$ window of vertex $i\le t$ is
$[2L-k-i,\ i+k)$ with $L=n-1$, whose endpoints satisfy $2L-k-i\ge L-k+1$ and
$i+k\le (L-1)+k$, so it is contained in the window of $t$ \emph{as an
interval}.  This is not enough, for two independent reasons.

First, the degrees carried by vertex $i$ obey the parity constraint
$d\equiv i+k\pmod2$, so the actual degree \emph{sets} need not be nested when
$i\not\equiv t\pmod 2$.  In $D_6$ with $\ka=3$ and $t=4$, vertex $3$ carries
multiplicity $2$ exactly at degree $4$ while vertex $t$ carries it exactly at
degrees $3,5$; the sets are disjoint.  (The same happens in $D_7$, $k=3$.)
Second, even granting containment, knowing that a repeat at $i$ occurs at a
degree in the $t$-window does not produce a repeat at $t$: nothing transports
$\dim V_i[d]\ge2$ to $\dim V_t[d']\ge2$.

The mechanism that does work is inward injectivity, Lemma \ref{lem:I} below,
which moves a repeat one vertex towards $t$ and shifts the degree by one --- and
with that shift the $D_6$ numbers do line up, $i=3$ at $d=4$ mapping to $t$ at
$d=5$.  Lemma \ref{lem:I} is stated and used in \S\ref{sec:typeE}; it is what
proves Theorem \ref{prop:lemA} in both types.
\end{remark}

\begin{proposition}[top truncation]\label{prop:T}
Let $N\subseteq I_\ka$ be a graded submodule with $\dimv N=v=\la-w\la$ for
some $w\in W$.  Then the colour-$t$ multiplicities of $N$ are the greedy
\emph{top} truncation of the sequence of Lemma \ref{lem:P} to total $v_t$:
the same degrees, filled from the highest downward.
\end{proposition}

Note that the sequence $(1,2^{k-1},1)$ of Lemma \ref{lem:P} is palindromic, so
its top truncations and its bottom truncations have the same multiset of
multiplicities, namely
\[
 v_t=0\mapsto(\,),\quad 1\mapsto(1),\quad 2\mapsto(1,1),\quad 3\mapsto(1,2),
 \quad 4\mapsto(1,2,1),\quad 5\mapsto(1,2,2),\ \dots
\]
Only the multiset is used below.

\begin{proof}[Proof for $k=2$, checked for $n\le7$]
Here the colour-$t$ profile of $I_\ka$ is $(1,2,1)$ at degrees $d_0,d_0+2,
d_0+4$ with $d_0=n-4$.  Computing the submodule generated by a single
colour-$t$ element shows that its colour-$t$ content is all of $(1,2,1)$ if the
element sits in degree $d_0$, is $(1,1)$ in degrees $d_0+2,d_0+4$ if it sits in
degree $d_0+2$ (for every choice of vector there), and is $(1)$ if it sits in
degree $d_0+4$.  Hence the only colour-$t$ profiles a submodule can have are
$(0,0,0),(0,0,1),(0,1,1),(0,2,1),(1,2,1)$, and each is a top truncation.  No
hypothesis on $v$ is needed.  The generation pattern was computed for
$D_4$--$D_7$ and is uniform in $n$; we have not written that induction out, but
nothing depends on it, since Theorem \ref{thm:staircase} gives Proposition
\ref{prop:T} for every $k$.
\end{proof}

\begin{remark}[how Proposition \ref{prop:T} is proved]\label{rem:Tk3}
Call a graded submodule $M\subseteq I_\ka$ a \emph{staircase} if its colour-$t$
multiplicities are the greedy top truncation of those of $I_\ka$ --- zero below
one degree, proper at it, full above --- and if at each of the two leaves $M$
occupies a final segment of degrees.  Everything left open in this paper is the
following single statement.

\begin{proposition}[the staircase property]\label{conj:staircase}
If $\dimv M$ is $w$-extremal then $M$ is a staircase.  Equivalently, by the
rigidity property of \cite{BKT} quoted in \S\ref{sec:setup}, a submodule that
is not a staircase satisfies $\delta(\dimv M)>0$, and so is never the unique
submodule of its dimension vector.
\end{proposition}

\noindent The proof occupies the rest of this section and is completed in
Theorem \ref{thm:staircase}: the statement is first made word-free, then reduced
to $k$ explicit weights, then to the closure of an explicitly written subspace,
and finally to an inequality between two integer sequences.

Proposition \ref{conj:staircase} gives Proposition \ref{prop:T} at once: an
extremal $v$ carries a unique submodule, namely $V=\soc_\sigma(I_\ka)$ by
Theorem \ref{thm:main}, and the staircase property applies to it.  It is
verified on all $1124$ graded submodules of $I_\ka$ over $D_4$--$D_7$: exactly
$78$ fail to be staircases and every one has $\delta=2$
(\texttt{master.py}).

\begin{proposition}\label{prop:tau}
Proposition \ref{conj:staircase} holds whenever $v_{n-1}=v_n$.
\end{proposition}

\begin{proof}
Suppose $M$ is not a staircase at $t$, say $M_t[d]\ne0$ while
$M_t[d+2]\subsetneq (I_\ka)_t[d+2]$.  Every nonzero $x\in M_t[d]$ is then
\emph{special} --- annihilated by an arrow from $t$ to a leaf --- since
otherwise the three double paths $a_{j\to t}a_{t\to j}x$, which satisfy one
linear relation and hence span at most two dimensions, would by direct
computation span all of $(I_\ka)_t[d+2]$, forcing it into $M$.  The special
vectors form two lines rather than a subspace, so $M_t[d]$ is one of them, and
$\dim (I_\ka)_t[d]=2$ at every such degree.  The diagram automorphism $\tau$
exchanging the leaves fixes $\ka$ (as $k\le n-2$), hence acts on $I_\ka$, and
exchanges the two arrows out of $t$ to the leaves, hence the two special lines;
so $\tau(M)\ne M$.  When $v_{n-1}=v_n$ we have $\dimv\tau(M)=\dimv M$, and two
distinct submodules of one dimension vector contradict rigidity, so $\dimv M$
is not extremal and $\delta>0$ by Corollary \ref{cor:pol}.
\end{proof}

The statement now reduces to a single assertion at a single vertex.  Carry
along the construction the hypothesis
\[
   (H)\qquad X \text{ is an up-set at } t \text{ and at the two leaves,}
\]
where \emph{up-set} means zero below one degree, proper at it, full above; at
$t$ this is top-truncation, and at a leaf, where $\dim (I_\ka)_j[d]\le1$, it
says $X_j$ occupies a final segment.  Steps at other colours do not change $X$
there, so $(H)$ need only be checked at colour-$t$ and leaf steps.

\emph{Leaf steps.}  Let $j$ be a leaf, whose only neighbour is $t$.  Then
$c_j=v'_t-2v'_j\ge1$ because the word is reduced.  The leaf degrees are
$g_1+1>\dots>g_k+1$ with $g_i=t+k-2i$, so under $(H)$ the largest unfilled one
is $D_j=t+k-2v'_j-1$, while $D=t+k-2a$ if $v'_t=2a$ and $D=t+k-2a-2$ if
$v'_t=2a+1$.  In the first case $c_j\ge1$ forces $v'_j\le a-1$, in the second
$v'_j\le a$; either way $D_j>D$.  As $X$ is an up-set at $t$ it is full above
$D$, hence at $D_j+1$, so the socle condition
$a_{j\to t}\xi\in X_j[D_j+1]$ is automatic and the top unfilled leaf slot lies
in $\soc_{(j)}(I_\ka/X)$ --- as does every unfilled leaf slot at a degree
$\ge D$, a consecutive block at the top.  So $(H)$ is preserved at leaf steps.

\emph{Colour-$t$ steps.}  Preservation at $t$ is (G1) and (G2), and both follow
from $(H)$, the bound $c\ge1$, and two inequalities about the profile of
$I_\ka$ alone:
\[
   \sum_{d\ge D+1}\dim (I_\ka)_{t-1}[d]\ \le\ 2a-[k=t]
   \quad\text{and}\quad
   \le\ 2a+2-[k=t]
\]
in the two cases.  These are statements about the profile of Lemma \ref{lem:P}
alone, and the strand form of it proves them by counting.  At $t-1$ the $+$
strand runs over $[\,|t-1-k|,\ t+k-3\,]$ and the $-$ strand over
$[\,t-k+3,\ 2t-|t-1-k|\,]$, in steps of $2$, each with $\min(t-1,k)$ slots.  In
the first case $D+1=t+k-2a+1$, which has the parity of $t-1+k$, so the $+$
strand contributes $\max(0,a-1)$ slots above it and the $-$ strand
$a+1-[k=t]$, capped at $\min(t-1,k)$; in the second $D+1=t+k-2a-1$ and the
contributions are $a$ and $a+2-[k=t]$.  Summing gives $2a-[k=t]$ and
$2a+2-[k=t]$, with equality in the generic range and the cap only helping.
(Checked at all $164160$ instances over $D_5$--$D_{80}$ and every $k$.)  At the leaves, $c\ge1$ gives
$\mu'_{n-1}\le0$, and then
$v'_{n-1}+v'_n=v'_t-\mu'_{n-1}$ with $|v'_{n-1}-v'_n|=|\mu'_n|\le1$ yields
$\min(v'_{n-1},v'_n)\ge a$ in the (G1) case, and in the (G2) case exactly one
leaf deficient by one dimension when $r'=0$ and none when $r'=2$.

\emph{What is left.}  Both (G1) and (G2) also need information about the chain
neighbour $t-1$ at one degree, and that is all that remains:
\[
   (*)\qquad
   \begin{array}{l}
   \text{at a colour-$t$ step, if } X_t[D]\ne0 \text{ then } X_{t-1}[D+1]
   \text{ is full,}\\
   \text{and if } X_t[D]=0 \text{ it is deficient by at most one dimension.}
   \end{array}
\]
This is strictly weaker than an up-set at $t-1$: $X_{t-1}$ is full in \emph{all}
degrees $>D$ at only $2207$ of $2337$ colour-$t$ steps, whereas $(*)$ holds at
all of them, so degree $D+1$ is exactly the right place to look.

Half of $(*)$ is proved.  Suppose $X_t[D]\ne0$, so it is spanned by a special
vector $x$, say $a_{t\to n}x=0$.  Then $a_{t\to(t-1)}x\ne0$: otherwise $x$ is
killed by two of the three arrows out of $t$, and not by the third, since
$I_\ka$ has simple socle $S_\ka$ and $\ka\le n-2$ is not a leaf; so
$y=a_{t\to(n-1)}x\ne0$, and the relation at $t$,
\[
   \epsilon_1 a_{(t-1)\to t}a_{t\to(t-1)}x+\epsilon_2 a_{(n-1)\to t}y
   +\epsilon_3 a_{n\to t}a_{t\to n}x=0,
\]
has vanishing first and third terms, giving $a_{(n-1)\to t}y=0$ --- a nonzero
element of the leaf $n-1$ killed by its only arrow, hence in
$\soc(I_\ka)=S_\ka$, forcing $\ka$ to be that leaf.  So
$a_{t\to(t-1)}(X_t[D])$ is a nonzero line in $X_{t-1}[D+1]$, and $(*)$ holds
whenever $\dim (I_\ka)_{t-1}[D+1]=1$ --- which is $422$ of the $547$ instances
with $X_t[D]\ne0$.

Splitting the $2337$ colour-$t$ steps by branch and by
$r'=|\mathrm{supp}\,\mu'\cap\{n-1,n\}|$ gives a clean table:
\[
\begin{array}{lcccc}
\text{branch} & n & \text{full} & \text{deficient by }1 & \text{by more}\\\hline
X_t[D]\ne0 & 547 & 547 & 0 & 0\\
X_t[D]=0,\ r'=0 & 855 & 855 & 0 & 0\\
X_t[D]=0,\ r'=1 & 935 & 805 & 130 & 0
\end{array}
\]
The first row is what (G1) requires and the other two what (G2) requires:
the deficiency never exceeds one and occurs only in the $r'=1$ branch, which is
precisely where both leaves are full, so at most one neighbour is ever
deficient.  What remains is therefore (A) the first row when
$\dim (I_\ka)_{t-1}[D+1]=2$, namely $125$ of the $547$; (B) the second row; and
(C) the third row, which is automatic unless that dimension is $2$.

All three concern the single space $X_{t-1}[D+1]$, and each would follow from
the up-set property at $t-1$, which is the statement
\[
   (*)'\qquad \text{at every colour-$t$ socle step, } X \text{ is an up-set at }
   t-1 ,
\]
verified at $2179/2179$ steps (\texttt{upsets.py}).  What makes $(*)'$
tractable is that the construction order can be removed from it entirely.  At
such a step $X=V(w')$ for a prefix $w'$ of the reduced word, so $X$ is itself
an extremal module by Theorem \ref{thm:main}; and the step having colour $t$
means $c_t=\langle\al_t^\vee,\mu'\rangle\ge1$ for $\mu'=\la-\dimv X$.  Hence
$(*)'$ follows from
\[
   (*)^{\circ}\qquad
   \begin{array}{l}
   \text{every extremal } M\subseteq I_\ka \text{ with }
   \langle\al_t^\vee,\la-\dimv M\rangle\ge1\\
   \text{is an up-set at } t-1 ,
   \end{array}
\]
which mentions neither $w$ nor any order of construction.  The side condition
is exactly right: over $D_5$--$D_{10}$ and all $\ka\le n-2$ in range,
$(*)^{\circ}$ holds at $1998/1998$ extremal modules, while among those failing
to be up-sets at $t-1$ --- there are $111$ --- \emph{every} one has
$c_t=0$ (\texttt{starchar.py}).

$(*)^{\circ}$ can now be reduced to finitely many weights, uniformly in the
rank.  Let $J=I\setminus\{t-1,t\}$ and let $j\in J$ be a descent of $\mu$.  Put
$\nu=s_j\mu$, which is shorter.  By Theorem \ref{thm:main}, $V(\mu)=\soc_j
V(\nu)$, and a socle step of colour $j$ enlarges the module only at vertex $j$;
as $j\ne t-1$,
\[
   V(\mu)\ \text{and}\ V(\nu)\ \text{have the same graded profile at } t-1 .
\]
So the property ``up-set at $t-1$'' is constant on $W_J$-orbits, and
$(*)^{\circ}$ for $\mu$ follows from $(*)^{\circ}$ for the $W_J$-minimal
representative of its orbit.  (One cannot take $j=t$: there
$\langle\al_t^\vee,s_t\mu\rangle=-\langle\al_t^\vee,\mu\rangle\le-1$, so $s_t\mu$
is longer.)

\begin{proposition}\label{prop:reps}
The $W_J$-minimal $\mu\in W\la$ with $\langle\al_t^\vee,\mu\rangle\ge1$ are
exactly those whose only descent is $t-1$, namely, in Bourbaki coordinates,
\[
   \mu_a=e_1+\dots+e_a-(e_{n-2-c}+\dots+e_{n-3})+e_{n-2},
   \qquad a+c=k-1,\quad 0\le a\le k-1 .
\]
There are $k$ of them, independently of $n$ --- or $k-1$ when $k=n-2$, where
$a=k-1$ gives $\mu=\la$.
\end{proposition}

\begin{proof}
For $k\le n-2$, $W\om_k=\{\sum_{i\in S}\eta_ie_i:|S|=k,\ \eta_i=\pm1\}$, so
$\mu=(\eta_1,\dots,\eta_n)$ with $\eta_i\in\{0,\pm1\}$ and exactly $k$ of them
nonzero, and $\langle\al_j^\vee,\mu\rangle=\eta_j-\eta_{j+1}$ for $j\le n-1$,
$\langle\al_n^\vee,\mu\rangle=\eta_{n-1}+\eta_n$.  The hypotheses read
$\eta_j\ge\eta_{j+1}$ for every $j\le n-1$ with $j\ne n-3$, together with
$\eta_{n-1}+\eta_n\ge0$ and $\eta_{n-2}-\eta_{n-1}\ge1$.  If $\eta_{n-1}=-1$
then $\eta_n\le-1$, contradicting $\eta_{n-1}+\eta_n\ge0$; if $\eta_{n-1}=1$
then $\eta_{n-2}\ge2$, impossible.  So $\eta_{n-1}=0$, whence $\eta_n\le0$ and
$\eta_n\ge0$, so $\eta_n=0$ and $\eta_{n-2}=1$.  On $1,\dots,n-3$ the sequence
is weakly decreasing, hence of the form $1^a0^b(-1)^c$ with $a+b+c=n-3$, and
counting nonzeros gives $a+c+1=k$.  Conversely each such $\mu_a$ has its unique
descent at $n-3=t-1$.  Finally $b=n-2-k\ge0$, so $a$ runs over $0,\dots,k-1$.
\end{proof}

The reduction is verified on the modules themselves in \texttt{reduce.py}:
$15364$ instances of $(*)^{\circ}$ over $D_4$--$D_{10}$, $k\le6$, collapse to
$83$ representatives, the profile at $t-1$ agreeing across every reduction step
($1705/1705$).

At the representatives more is true than $(*)^{\circ}$ asks: $V(\mu_a)$ is the
greedy top truncation of $I_\ka$ at \emph{every} vertex, not only at $t$ where
Proposition \ref{prop:T} puts it.  So what remains is
\[
   (B)\qquad
   \begin{array}{l}
   \text{for the $k$ weights $\mu_a$ of Proposition \ref{prop:reps}, the graded}\\
   \text{multiplicities of $V(\mu_a)$ are the greedy top truncation of those}\\
   \text{of $I_\ka$, at every vertex,}
   \end{array}
\]
which gives $(*)^{\circ}$ at once, a greedy top truncation being an up-set.
$(B)$ holds at every representative in computable range ($43/43$).  Its point
is that it is a bounded list: the same assertion for arbitrary extremal $\mu$
is false --- see below --- and even for arbitrary leaf-balanced extremal $\mu$
with $c_t\ge1$ it fails, at $62$ of $348$ instances over $D_4$--$D_8$.  The
earlier route through a chainwide staircase invariant (\texttt{staircase.py})
required in addition control of how colour-$(t-j)$ steps interleave with
colour-$t$ steps along the $c$-trace; nothing of the kind is needed here.

$(B)$ can be turned into a closure check on an explicitly written subspace.  A
graded top truncation is not determined by its multiplicities: where $I_\ka$
carries multiplicity $2$ and the truncation takes only $1$, a line must be
chosen.  Over all extremal modules exactly three lines ever occur in such a
slot, and the representatives always use the same one, canonically described:
$\tau$, the diagram automorphism exchanging the leaves, fixes $I_\ka$ (as $\ka$
is a chain vertex), acts on every multiplicity-$2$ slot with eigenvalues
$\pm1$, and the line used is the $+1$ eigenline; the other two are the special
lines, the kernels of the two leaf arrows, which $\tau$ exchanges.  Define
$T(v)$, the \emph{canonical top truncation}, to be full above the threshold,
zero below, and the $\tau$-eigenline at a partially taken multiplicity-$2$
slot.  Then $\dimv T(v)=v$ by construction, so by rigidity
\[
   (B)\quad\Longleftrightarrow\quad T(v_a)\ \text{is a submodule of } I_\ka ,
\]
and the socle construction disappears from the statement altogether.

Closure is then almost formal, given three facts about $I_\ka$ and $\tau$
alone --- no words, no dimension vectors:
\begin{itemize}\itemsep0pt
\item[(T0)] every arrow between graded slots has rank
$\min(\dim\text{source},\dim\text{target})$: it is injective or surjective,
whichever is possible ($450/450$);
\item[(T1)] a rank-$1$ arrow into a multiplicity-$2$ chain slot has image the
$\tau$-eigenline ($22/22$);
\item[(T2)] the $\tau$-eigenline of a multiplicity-$2$ chain slot is
annihilated by exactly the arrows into $1$-dimensional \emph{chain} slots ---
zero there ($22/22$), nonzero into $2$-dimensional chain slots ($36/36$) and
into both leaves ($44/44$).
\end{itemize}
Every closure obligation for $T(v_a)$ is then of a resolved kind: full into
full (nothing to check); full into partial (the arrow has rank $1$ by (T0) and
its image is the eigenline by (T1)); eigenline into full; eigenline into
partial (the image is $\tau$-fixed, hence in the target eigenline, by
equivariance); and eigenline into an empty slot (the image vanishes by (T2)).
Whether an obligation is of one of these kinds depends only on the two slot
dimensions, on whether the target is a leaf, and on the greedy thresholds --- so
all of $(B)$ beyond (T0)--(T2) is arithmetic.

These three have a common source, and with it they stop being conjectural.  At a
chain vertex $\tau$ splits every graded piece into eigenspaces --- write
$a^{\pm}_d(j)$ for the eigen-multiplicities --- and that splitting is as clean
as it could be.  With $t=n-2$, so that the top degree is $2t$:
\begin{itemize}\itemsep0pt
\item[(W1)] each eigenspace is multiplicity-free in every chain slot, so a
multiplicity-$2$ slot is exactly one $+$ line and one $-$ line ($8255/8255$);
\item[(W2)] in closed form,
\[
   a^{+}_d(j)=1 \iff |j-k|\le d\le j+k-2 \ \text{ and }\ d\equiv j+k\ (2),
   \qquad a^{-}_d(j)=a^{+}_{2t-d}(j),
\]
the $+$ strand being the Dirichlet half-line wave of the generator and the $-$
strand its reflection in the top degree.  The two add up to the multiplicity of
Lemma \ref{lem:P} ($27430/27430$) and are the eigen-multiplicities ($8255/8255$);
\item[(W3)] hence each strand occupies a contiguous run of degrees at every
chain vertex, the $+$ strand the low ones and the $-$ strand the high ones,
overlapping exactly in the multiplicity-$2$ window ($2236/2236$);
\item[(W4)] every arrow restricts to an \emph{isomorphism} between consecutive
nonzero slots of the same strand ($17706/17706$).
\end{itemize}
So $I_\ka$ is two multiplicity-free waves travelling through the chain, the
arrows transporting each (\texttt{strands.py}).

Of these, (W1) and (W3) are immediate from the closed form (W2), and (W2) itself
is a consequence of Lemma \ref{lem:R}.

\begin{proposition}\label{prop:W2}
(W2) holds; hence so do (W1) and (W3).
\end{proposition}

\begin{proof}
The complex of Lemma \ref{lem:R} at a chain vertex $j$,
$(I_\ka)_j[d-1]\to\bigoplus_{i\sim j}(I_\ka)_i[d]\to(I_\ka)_j[d+1]$, consists of
$\tau$-equivariant maps: the arrows between chain vertices commute with $\tau$,
and for $j=t$ the two leaf summands are exchanged by it, so their sum is
$\tau$-stable.  As $\operatorname{char}\ne2$ the complex splits into
eigen-parts, and exactness and the injectivity of the first map below the top
degree are inherited by each summand.  Writing $b_d$ for the common dimension of
$(I_\ka)_{n-1}[d]$ and $(I_\ka)_n[d]$, the pair is the induced representation of
$\mathbb{Z}/2$, so each eigen-part of it has dimension $b_d$.  Counting
dimensions in the $\epsilon$-part therefore gives, for every chain vertex $j$ and
$0<d<2t$,
\[
   a^{\epsilon}_{d+1}(j)+a^{\epsilon}_{d-1}(j)
   =\sum_{i\sim j,\ i\le t}a^{\epsilon}_d(i)+[\,j=t\,]\,b_d ,
\]
while Lemma \ref{lem:R} at a leaf reads
$b_{d+1}+b_{d-1}=a^{+}_d(t)+a^{-}_d(t)$.  The generator $e_k$ spans
$(I_\ka)_k[0]$ and is fixed by $\tau$, so $a^{+}_0(j)=\delta_{jk}$,
$a^{-}_0\equiv0$ and $b_0=0$.  Each equation expresses degree $d+1$ in terms of
degrees $d$ and $d-1$, so this system has at most one solution, determined by
induction on $d$; and the closed form of (W2), with $b_d=a^{+}_d(t+1)$, satisfies
it.  (That last verification is arithmetic in $n,k,j,d$, and holds over
$D_5$--$D_{60}$ for every $k$ at all $11\,443\,152$ instances of the chain
recursion and $130\,032$ of the leaf recursion, with no exception.)
\end{proof}

(W4) is a consequence of the same $\tau$-refined Lemma \ref{lem:R}, through a
propagation along the chain.  Write $[\mathrm{lo}^{\epsilon}_j,\
\mathrm{hi}^{\epsilon}_j]$ for the interval (W3) attaches to the strand
$\epsilon$ at the chain vertex $j$, so that
$\mathrm{hi}^{+}_j=j+k-2$ and $\mathrm{hi}^{-}_j=2t-|j-k|$.  The single
arithmetic fact used is
\begin{equation}\label{eq:hi}
   \min\bigl(\mathrm{hi}^{\epsilon}_{j-1},\mathrm{hi}^{\epsilon}_{j+1}\bigr)
   =\mathrm{hi}^{\epsilon}_j-1 ,
\end{equation}
immediate for $+$, and for $-$ because
$\max(|j-1-k|,|j+1-k|)=|j-k|+1$.

\begin{lemma}[a death propagates]\label{lem:death}
Let $1<j<t$ and let $x$ span the $\epsilon$-line at $(d,j)$.  Suppose
$a_{j\to i}x=0$ for one chain neighbour $i$ of $j$, the slot $(d+1,i)$ carrying
$\epsilon$.  Write $i'$ for the other chain neighbour.  Then $a_{j\to i'}x$
spans the $\epsilon$-line at $(d+1,i')$, the slot $(d+2,j)$ carries $\epsilon$,
and $a_{i'\to j}\bigl(a_{j\to i'}x\bigr)=0$.
\end{lemma}

\begin{proof}
As $j<t$ its only neighbours are $i$ and $i'$, so the $\epsilon$-part of Lemma
\ref{lem:R} forces $a_{j\to i'}x\ne0$; it lies in the $\epsilon$-part of
$(I_\ka)_{i'}[d+1]$, which is at most a line by (W1), so that slot carries
$\epsilon$ too.  Both slots $(d+1,j\pm1)$ therefore carry $\epsilon$, whence
$d+1\le\min(\mathrm{hi}^{\epsilon}_{j-1},\mathrm{hi}^{\epsilon}_{j+1})
=\mathrm{hi}^{\epsilon}_j-1$ by \eqref{eq:hi}, i.e.\
$d+2\le\mathrm{hi}^{\epsilon}_j$; and $d+2\ge d\ge\mathrm{lo}^{\epsilon}_j$ with
the parity unchanged, so $(d+2,j)$ carries $\epsilon$.  Finally the relation at
$j$ reads $\epsilon_1a_{i\to j}a_{j\to i}x+\epsilon_2a_{i'\to j}a_{j\to i'}x=0$,
whose first term vanishes.
\end{proof}

\begin{proposition}[outward arrows]\label{prop:W4out}
For either strand, $a_{j\to(j+1)}$ is nonzero on the $\epsilon$-line at $(d,j)$
whenever $(d+1,j+1)$ carries $\epsilon$.
\end{proposition}

\begin{proof}
Suppose $a_{j\to(j+1)}x=0$.  If $j=1$ the only chain neighbour is $2$ and the
$\epsilon$-part of Lemma \ref{lem:R} is contradicted at once.  Otherwise Lemma
\ref{lem:death} with $i=j+1$, $i'=j-1$ reproduces the situation at $(d+1,j-1)$:
its outward arrow dies and the target $(d+2,j)$ carries $\epsilon$.  Iterating
lowers the vertex to $1$.
\end{proof}

\begin{proposition}[the folded leaf]\label{prop:W4l}
The map $J[d]^{\epsilon}\to(I_\ka)_t[d+1]^{\epsilon}$,
$J=(I_\ka)_{n-1}\oplus(I_\ka)_n$, is nonzero whenever both sides are.
\end{proposition}

\begin{proof}
The leaves are multiplicity-free, so $(I_\ka)_{n-1}[d]$ is a line, spanned by
$u$ say, and $J^{\epsilon}$ is spanned by $(u,\epsilon\tau u)$, whose image is
$(1+\epsilon\tau)z$ with $z=a_{(n-1)\to t}u$.  As $n-1$ has the single
neighbour $t$, Lemma \ref{lem:R} makes $z\ne0$, and the relation at $n-1$ reads
$a_{t\to(n-1)}z=0$.  If $\dim(I_\ka)_t[d+1]=1$ the claim is exactly $z\ne0$.

So let that slot be $2$-dimensional and suppose $z$ were a $\tau$-eigenvector,
of eigenvalue $\epsilon$.  Applying $\tau$ to $a_{t\to(n-1)}z=0$ gives
$a_{t\to n}z=0$, so both leaf arrows kill $z$; by Lemma \ref{lem:R} at $t$,
$y:=a_{t\to(t-1)}z\ne0$, and $y$ lies in the $\epsilon$-part at $(d+2,t-1)$,
which is therefore nonzero.  The relation at $t$ applied to $z$ has both leaf
terms zero, so $a_{(t-1)\to t}y=0$.  Finally $(d+3,t)$ carries $\epsilon$: for
$\epsilon=+$ because $d+2\le\mathrm{hi}^{+}_{t-1}=\mathrm{hi}^{+}_t-1$; for
$\epsilon=-$ because the slot $(d+1,t)$ carries $+$ as well, so
$d+1\le\mathrm{hi}^{+}_t=t+k-2$ and hence $d+3\le t+k=\mathrm{hi}^{-}_t$.  Thus
the outward arrow at $(d+2,t-1)$ dies with its target carrying $\epsilon$,
contradicting Proposition \ref{prop:W4out}.  So $z$ is not an eigenvector and
both of its eigen-components are nonzero.
\end{proof}

\begin{corollary}\label{cor:W4}
(W4) holds.
\end{corollary}

\begin{proof}
Outward arrows are Proposition \ref{prop:W4out}.  For an inward arrow, a death
at $(d,j)$ towards $j-1$ produces by Lemma \ref{lem:death} a death at
$(d+1,j+1)$ towards $j$, again inward, and iterating raises the vertex to $t$.
There the relation carries the two leaf terms: the death makes the composite
through the leaves vanish on the $\epsilon$-line, while Lemma \ref{lem:R} makes
its image in $J^{\epsilon}$ nonzero --- contradicting Proposition
\ref{prop:W4l}.
\end{proof}

With (W1)--(W4) in place, (T0), (T1) and (T2) are unconditional.

\begin{proposition}\label{prop:T012}
(W) implies (T0), (T1) and (T2) for arrows between chain vertices.
\end{proposition}

\begin{proof}
By (W4) the rank of an arrow $(d,j)\to(d+1,i)$ is the number of strands carried
by both slots, so (T0) says exactly that the two sets of strands present are
\emph{nested}.  Suppose not; then one slot carries only $+$ and the other only
$-$.  Throughout $i,j,k\le t$ and $|i-j|=1$.

\emph{Case A: $+$ only at $(d,j)$, $-$ only at $(d+1,i)$.}  Since $+$ holds at
$(d,j)$ we have $d\ge|j-k|\ge|i-k|-1$, so $d+1\ge|i-k|$ and the failure of $+$
at $(d+1,i)$ must be $d+1>i+k-2$, i.e.\ $d\ge i+k-2$; while $+$ at $(d,j)$ gives
$d\le j+k-2$.  If $i=j+1$ this reads $d\le i+k-3<i+k-2\le d$.  If $i=j-1$ then
$d\in\{i+k-2,i+k-1\}$ and the parity $d\equiv j+k=i+k+1$ forces $d=i+k-1=j+k-2$;
now $-$ at $(d+1,i)=(i+k,i)$ needs $i+k\ge 2t-i-k+2$, i.e.\ $i+k\ge t+1$, while
the absence of $-$ at $(d,j)$ needs $j+k-2<2t-j-k+2$, i.e.\ $j+k\le t+1$ ---
impossible for $j=i+1$.

\emph{Case B: $-$ only at $(d,j)$, $+$ only at $(d+1,i)$.}  From $+$ at
$(d+1,i)$, $d\le i+k-3$.  The absence of $+$ at $(d,j)$ is either $d>j+k-2$,
which with $d\le i+k-3$ forces $j<i-1$, or $d<|j-k|$; in the latter case $-$ at
$(d,j)$ gives $2t-j-k+2\le d<|j-k|$, which reads $2t+2<2j$ if $j\ge k$ and
$2t+2<2k$ if $j<k$, both impossible.

For (T1) the source carries one strand and the target two; suppose the source
strand were $-$.  From $+$ at $(d+1,i)$, $d\le i+k-3$.  The absence of $+$ at
$(d,j)$ is $d>j+k-2$, contradicted by $d\le i+k-3\le j+k-2$, or $d<|j-k|$,
excluded exactly as in Case B.

For (T2) the source carries both strands and the target one; suppose the target
strand were $+$.  The absence of $-$ at $(d+1,i)$ is either $d\le 2t-i-k$,
which against $d\ge 2t-j-k+2$ from $-$ at $(d,j)$ gives $i\le j-2$; or
$d\ge 2t-|i-k|$, which against $d\le j+k-2$ from $+$ at $(d,j)$ gives
$2t+2\le i+j$ when $i\ge k$ and $2t+2\le j-i+2k\le 2t+1$ when $i<k$.  All are
impossible.
\end{proof}

The nesting statement, and (T1) and (T2) with it, are checked directly from the
closed form over $D_5$--$D_{100}$, every $k$, every adjacent pair and every
degree: $23\,517\,840$ nesting instances and $156\,848$ each of (T1) and (T2),
with no exception (\texttt{strands.py}).

In that form everything can be tested far past the range where the modules can
be built, and the residue reduces to two integer sequences.  With (W) proved, a
$\tau$-invariant graded submodule splits at chain vertices as
$N^{+}\oplus N^{-}$ with $N^{\epsilon}$ a subobject of the strand $\epsilon$;
and since $(d,j)$ reaches $(d+2,j)$ by an outward step followed by an inward
one, a subobject meets each vertex in a \emph{top interval}.  So such a
submodule is determined by the numbers $\dim N^{\epsilon}_j$, and closure of
$T(v_a)$ says exactly that its two strand-parts are up-closed.

Each strand has $M_j=\min(j,k)$ slots at a chain vertex $j$, of which
$o_j=\max(0,j+k-t-1)$ are shared with the other ($26486/26486$).  Greedy from
the top fills the $M_j-o_j$ slots carrying only the $-$ strand, then the $o_j$
shared ones two dimensions at a time --- an odd one taking the $\tau$-line,
which is the $+$ strand --- and then the $+$-only slots, so the occupation
numbers are
\[
   (s^{+}_j,s^{-}_j)=
   \begin{cases}
     (0,\,v_j), & v_j\le M_j-o_j,\\[2pt]
     \bigl(\lceil r/2\rceil,\ M_j-o_j+\lfloor r/2\rfloor\bigr),
       & r:=v_j-(M_j-o_j)\le 2o_j,\\[2pt]
     (v_j-M_j,\ M_j), & \text{otherwise,}
   \end{cases}
\]
which reproduces the truncation exactly ($459282/459282$).  Writing
$\mathrm{hi}^{+}_j=j+k-2$ and $\mathrm{hi}^{-}_j=2t-|j-k|$ for the top of each
strand, the lowest occupied degree is
$\mathrm{hi}^{\epsilon}_j-2(s^{\epsilon}_j-1)$, so up-closedness at an edge
$j\to j+1$ is precisely
\[
   0\le s^{+}_{j+1}-s^{+}_j\le1,
   \qquad |s^{-}_{j+1}-s^{-}_j|\le1 .
\]
Write $u_j=v_j-\min(j,k)$, so that
$s^{+}_j=\max\bigl(0,\ u_j,\ \lceil(u_j+o_j)/2\rceil\bigr)$
($459282/459282$) and $s^{-}_j=v_j-s^{+}_j$.  Almost all of the two inequalities
now falls out.

\begin{proposition}\label{prop:incr}
$0\le s^{+}_{j+1}-s^{+}_j\le1$ for every chain edge; and
$|s^{-}_{j+1}-s^{-}_j|\le1$ except possibly where $v$ jumps by $2$.
\end{proposition}

\begin{proof}
Reading off $x_p=\la_p-(\mu_a)_p$ region by region gives
$u_{j+1}-u_j\in\{-1,0,1\}$ always, and $=-1$ only for $j+1\le a$ or $j+1=t$;
also $o_{j+1}-o_j\in\{0,1\}$.  Each of the three terms in $s^{+}$ increases by
at most one under such a step, so $s^{+}_{j+1}-s^{+}_j\le1$.  For the lower
bound, only the two exceptional steps can lower a term.  For $j+1\le a$ one has
$v_j=0$ and $u_j+o_j\le k-t-1<0$, so $s^{+}$ is $0$ at both ends.  For $j+1=t$,
$o$ steps up while $u$ steps down, so $u_j+o_j$ is unchanged and only the
middle term $u$ falls; but $\lceil(u+o)/2\rceil<u$ forces $u\ge o+2$, whereas
$u_{t-1}\le o_{t-1}+1$ ($34104/34104$), so the maximum is not attained by $u$
alone and $s^{+}$ does not drop.  Finally $s^{-}_{j+1}-s^{-}_j=x_{j+1}-
(s^{+}_{j+1}-s^{+}_j)$ with $x_{j+1}\in\{-1,0,1,2\}$; the values $0$ and $1$ are
harmless, and $x_{j+1}=-1$ occurs only at $j+1=t$, where $s^{+}$ is unchanged by
the above.
\end{proof}

The one implication left over is proved by the same bookkeeping.

\begin{lemma}\label{lem:jump2}
If $x_{j+1}=2$ then $s^{+}_{j+1}-s^{+}_j=1$.
\end{lemma}

\begin{proof}
$x_p=2$ exactly for $a+b<p\le\min(k,t-1)$, where $b=t-k$.  So $j=p-1$ satisfies
$a+b\le j\le\min(k,t-1)-1$; in particular $j<k$, so $\min(j,k)=j$,
$\min(j+1,k)=j+1$ and $u_{j+1}-u_j=x_{j+1}-1=1$.  In this range
$v_j=b+2(j-a-b)$, hence $u_j=j-2a-b$ and $o_j=\max(0,j-b-1)$.

If $j\ge b+1$ then $o_j=j-b-1$ and $o_{j+1}=j-b$, so $u+o$ increases by $2$ and
$\lceil(u+o)/2\rceil$ by exactly $1$.  Moreover $u_j-o_j=1-2a\le1$, so
$u_j\le o_j+1$ and the middle term does not strictly beat
$\lceil(u_j+o_j)/2\rceil$; and $u_j+o_j=2(j-a-b)-1\ge-1$, so that term is
$\ge0$.  Hence $s^{+}_j=\lceil(u_j+o_j)/2\rceil$ and
$s^{+}_{j+1}\ge s^{+}_j+1$, which with Proposition \ref{prop:incr} is equality.

If $j\le b$ then $a+b\le j\le b$ forces $a=0$ and $j=b$, so $u_j=0$ and
$o_j=o_{j+1}=0$, giving $s^{+}_j=0$ and $u_{j+1}=1$, whence $s^{+}_{j+1}=1$.
\end{proof}

The obligations that involve a leaf are settled by three identities at the node.
At $t$ one has $v_t=2k-2a-2$ and each leaf carries $v=k-a-1$ ($34104/34104$);
$\min(t,k)=k$ and $o_t=k-1$, so the truncation at $t$ takes $k-a$ slots and its
lowest taken degree, which is also its threshold $D_t$, equals $t-k+2a+2$
($32508/32508$); while the leaves, being multiplicity-free with top degree
$t+k-1$, start at $t-k+2a+3$ ($32508/32508$) --- exactly one higher.  So every
taken slot at $t$ maps into a taken leaf slot, and every taken leaf slot maps
into the slot of degree $D_t+2$ at $t$, which is full.  Both leaf patterns are
therefore ``target full'', and require nothing ($1462902/1462902$ instances).

\begin{theorem}\label{thm:staircase}
Proposition \ref{conj:staircase} holds, and with it Theorem \ref{thm:typeD}.
\end{theorem}

\begin{proof}
By Proposition \ref{prop:reps} it suffices to prove $(B)$ at the $k$
representatives, and by rigidity that is the assertion that the canonical
truncation $T(v_a)$ is a submodule.  Its closure obligations along the chain are
the up-closedness of its two strand-parts, which by Proposition \ref{prop:incr}
and Lemma \ref{lem:jump2} is the pair of increment bounds; those involving a
leaf land in full slots by the identities above.  (T0), (T1) and (T2), which
identify the images, are Proposition \ref{prop:T012}.
\end{proof}

Every step above is re-checked in \texttt{audit.py}, which walks the chain from
the main theorem to the closed form and prints one line per link.  The counts,
over $D_5$--$D_{45}$ with every $k$ and every $a$ unless stated otherwise:
$\dimv\soc_\sigma(I_\ka)=v$ and agreement with the independent solver on $277$
elements; Lemma \ref{lem:I} at $315$ triples in type $E$ (the complete finite
verification) and $55$ in type $D$; Lemma \ref{lem:P} at $880$ profiles;
$v_t=r+2m$ at $24312$ elements; (W1)--(W4) at $3570$, $3570$, $1140$, $7300$
slots and the folded-leaf arrow at $635$; (T0), (T1), (T2) at $2\,922\,500$,
$32508$, $32508$ instances over $D_5$--$D_{60}$; the representatives at $24$
pairs $(n,k)$; $T(v_a)$ a submodule and equal to $V(\mu_a)$ at $32$; the
arithmetic at $8\,047\,275$ up-closure instances, $445\,178$ increments of each
kind, $148\,994$ drops, $84273$ jumps of two and $446\,982$ leaf obligations,
with the parity case occurring $0$ times; (G1) and (G2) at $164\,160$ instances
over $D_5$--$D_{80}$; and the criterion itself at $5819$ elements over
$D_4$--$D_{10}$.  The audit also runs negative controls --- statements that must
\emph{fail}, including the one-neighbour interval step that an earlier draft of
Lemma \ref{lem:death} wrongly used --- checks that the criterion is unchanged at
$p=5,7,11,13$ and $1000003$, and that it does not depend on the reduced word.

Extremality cannot be dropped from $(*)^{\circ}$.  The candidate localisation
--- for every graded $N\subseteq I_\ka$, $(H)$ and $c_t\ge1$ imply $N$ is an
up-set at $t-1$ --- holds for every graded submodule with $\ka\le3$ and fails
at $\ka=4$, where $t-1$ first carries multiplicity $2$ in two different
degrees.  The smallest counterexample is $D_6$, $\ka=4$, $t=4$, with graded
profile ($\dim N/\dim I_\ka$ by degree)
\[
\begin{array}{ll}
 1:\ 0/1\ \ 1/1 &
 4:\ 0/1\ \ 0/2\ \ 1/2\ \ 2/2\ \ 1/1\\
 2:\ 0/1\ \ 1/2\ \ 1/1 &
 5:\ 0/1\ \ 1/1\ \ 1/1\ \ 1/1\\
 3:\ 0/1\ \ 1/2\ \ 1/2\ \ 1/1 &
 6:\ 0/1\ \ 0/1\ \ 1/1\ \ 1/1
\end{array}
\]
which satisfies $(H)$, has $c_t=1$, and is partial at two degrees at $t-1$
(\texttt{localstar.py}).

It is worth recording what does \emph{not} work, since it is the first thing
one tries.  The hypothesis that $V$ is simply the greedy top truncation of
$I_\ka$ at \emph{every} vertex --- which would be a closed formula and would
settle everything at once --- is false.  It holds at $1185$ of $1307$ extremal
modules over $D_4$--$D_7$, $E_6$, $E_7$ and $A_5$, the smallest failure being
$D_5$, $\la=\om_2$, $\mu=(1,-1,0,0,0)$, where $V_2$ occupies degrees $2$ and
$6$ while the top truncation would give $6$ and $4$; vertex $2$ is the chain
neighbour $t-1$.  That element has $r+2m=2$ and so does have simple spectrum,
so the failure is not confined to degenerate cases.  Top truncation at $t$ and
at the two leaves --- the hypothesis $(H)$ --- does hold for every extremal
module tested.

Two facts delimit the argument.  It genuinely needs extremality: the
staircase property fails for $78$ of the $1124$ graded submodules above, and
the analogous statement for arbitrary submodules is false.  And the obvious
mechanism is not available: when $X_t[D]$ is a line it is \emph{always} one of
the two special lines ($527/527$), so one leaf arrow annihilates it and that
leaf's fullness at $D+1$ comes from elsewhere in $X$.
\end{remark}

\begin{proof}[Proof of Theorem \ref{thm:typeD}]
For the spin weights $\la$ is minuscule and Proposition \ref{prop:minuscule}
applies.  Let $k\le n-2$.  By Lemma \ref{lem:P} the colour-$t$ profile of
$I_\ka$ is $(1,2^{k-1},1)$.  By Theorem \ref{thm:main}, $V$ is a submodule of
$I_\ka$ with $\dimv V=v$, so Proposition \ref{prop:T} applies to it and the
colour-$t$ multiplicities of $V$ are the truncation displayed above.  Every
multiplicity is $\le1$ exactly when $v_t\le2$.  By Theorem \ref{prop:lemA},
simplicity at $t$ implies simplicity everywhere.  Lemma
\ref{lem:C} identifies $v_t=r+2m$.
\end{proof}

\begin{remark}
The disjunctive form ``$c_{\max}=1$ or $v_t\le2$'' of the criterion, natural
from the $c$-trace, collapses once the spin weights are separated: for
non-spin $\la$ in type $D$, $c_{\max}=1$ forces
$(r,m)\in\{(0,0),(0,1),(1,0),(2,0)\}$, all of which satisfy $r+2m\le2$.
\end{remark}

\section{Type \texorpdfstring{$E$}{E}}\label{sec:typeE}

Theorem \ref{thm:main} and Corollary \ref{cor:crit} are uniform and therefore
settle type $E$ as well: the criterion is decidable from any reduced word.
Two further points deserve comment.

\subsection{Reduction to the trivalent node}
As Remark \ref{rem:nowindow} explains, no containment-of-windows argument
proves Theorem \ref{prop:lemA}, in either type: in $E_6$ with $\ka=3$ (so
$I_3=\Pia e_5$) vertex $3$ carries multiplicity $2$ in path-length degree $4$
while vertex $4$ is zero there, and the same phenomenon already occurs in
$D_6$ with $\ka=3$.  The mechanism that does work is the following.

For $i\ne t$ let $\mathrm{nx}(i)$ denote the neighbour of $i$ one step closer
to $t$.

\begin{lemma}[inward injectivity]\label{lem:I}
Let $\Gamma$ be $E_6$, $E_7$ or $E_8$.  For every fundamental weight $\ka$,
every $i\ne t$ and every degree $d$ with $\dim (I_\ka)_i[d]\ge2$, the arrow
\[
   (I_\ka)_i[d]\longrightarrow (I_\ka)_{\mathrm{nx}(i)}[d+1]
\]
is injective.
\end{lemma}

\begin{proof}
The three diagrams have finitely many triples $(\ka,i,d)$ in total, so this is
a finite verification: $11$ instances in $E_6$, $55$ in $E_7$ and $249$ in
$E_8$, all injective.  See \texttt{inwardlemma.py}.
\end{proof}

\begin{proof}[Proof of Theorem \ref{prop:lemA} in type $E$]
Suppose $\dim V_i[d]\ge2$ for some $i\ne t$.  Then
$\dim (I_\ka)_i[d]\ge \dim V_i[d]\ge2$, so by Lemma \ref{lem:I} the inward
arrow is injective on $(I_\ka)_i[d]$, hence on the subspace $V_i[d]$.  Since
$V$ is a submodule of $I_\ka$, the image lies in $V_{\mathrm{nx}(i)}[d+1]$, so
$\dim V_{\mathrm{nx}(i)}[d+1]\ge2$.  Iterating $\mathrm{dist}(i,t)$ times gives
$\dim V_t[d+\mathrm{dist}(i,t)]\ge2$.  Contrapositively, if $V_t$ is
multiplicity-free then so is $V$.
\end{proof}

In type $D$ the same statement holds, and there it can be proved outright,
uniformly in the rank.

\begin{lemma}[Lemma \ref{lem:I} in type $D$]\label{lem:Id}
Let $\Gamma=D_n$, $t=n-2$, let $i<t$ be a chain vertex and let $d$ satisfy
$\dim (I_\ka)_i[d]\ge2$.  Then the inward arrow
$(I_\ka)_i[d]\to(I_\ka)_{i+1}[d+1]$ is injective.
\end{lemma}

\begin{proof}
Suppose $x\ne0$ lies in the kernel.  For $i<t$ the vertex $i$ has exactly the
two neighbours $i-1$ and $i+1$.  By Lemma \ref{lem:R} the socle of $I_\ka$ is
simple and concentrated in the top degree, so an element killed by every arrow
out of $i$ is either zero or that socle generator; since
$\dim (I_\ka)_i[d]\ge2$ we are below the top degree, and therefore
\[
   y:=a_{i\to i-1}(x)\ne0 .
\]
The preprojective relation at $i$ reads
$\varepsilon\, a_{i+1\to i}a_{i\to i+1}+\varepsilon'\,
a_{i-1\to i}a_{i\to i-1}=0$ on $(I_\ka)_i[d]$; applied to $x$, whose first
term vanishes by hypothesis, it gives $a_{i-1\to i}(y)=0$.  So $y$ is a nonzero
element of $(I_\ka)_{i-1}[d+1]$ killed by \emph{its} inward arrow: the same
situation, one vertex further from $t$ and one degree higher.

Iterating $i-1$ times produces a nonzero $z\in (I_\ka)_1[d+i-1]$ killed by the
arrow $1\to2$.  In $D_n$ the vertex $1$ has $2$ as its only neighbour, so $z$
is killed by every arrow out of $1$ and hence lies in $\soc(I_\ka)=S_\ka$,
forcing $\ka=1$.  But by Lemma \ref{lem:P} a chain vertex $i$ carries
multiplicity $2$ in exactly $\max(0,i+k-n+1)$ degrees, which vanishes for
$k=1$ and $i\le n-2$; so $\dim (I_\ka)_i[d]\ge2$ is impossible when $\ka=1$.
Contradiction.
\end{proof}

\begin{remark}\label{rem:Idtype}
The two leaves $n-1,n$ need no argument: by the folding in Lemma \ref{lem:P}
they satisfy $a_d(n-1)=a_d(n)=b_d$ with $2b_d=a_d(t+1)\le2$, so $b_d\le1$ and
Lemma \ref{lem:Id} is vacuous there.  With Lemma \ref{lem:Id} in hand the
proof of Theorem \ref{prop:lemA} given above applies verbatim in type $D$, for
every rank.  It is also confirmed directly for $D_4$--$D_{10}$ ($126$
instances) by \texttt{inwardlemma.py}, and the two inputs --- simple socle and
the relation at a chain vertex --- by \texttt{inwardproof.py}.
\end{remark}

\subsection{No closed form of type-$D$ shape}
What does \emph{not} carry over is the closed form of Theorem \ref{thm:typeD}:
$v_t$ is no longer a complete invariant.  In $E_6$ with $\la=\om_3$ both simple
and non-simple elements occur at $v_t=3$; likewise in $E_7$ with $\la=\om_2$ at
$v_t=4$.  A closed form must therefore involve more than the multiplicity at
the trivalent node.  We have not found one.

The colour-$t$ multiplicity profiles of $I_\ka$ in type $E$ are finite data,
computed once and for all; in $E_6$ they are $(1,1,1,1)$, $(1,1,2,1,1)$,
$(1,2,2,2,1)$, $(1,2,3,3,2,1)$, $(1,2,2,2,1)$, $(1,1,1,1)$ for
$k=1,\dots,6$, and the maximum multiplicity anywhere in $I_\ka$ is $6$,
attained in $E_8$ at $k=4$.  This is also why the truncation mechanism of
\S\ref{sec:typeD} is specific to type $D$: it turns on
$\dim (I_\ka)_t[d]\le2$, so that the three double paths at $t$, subject to one
relation, can span everything.

\section{Verification}\label{sec:verification}

The accompanying code reproduces every claim; each script runs standalone and
prints a pass/fail line.

\begin{itemize}
\item \texttt{proof.py} --- the five steps of \S\ref{sec:proof}.  The
Crawley-Boevey identity holds at $5517/5517$ socle steps; $\dimv
I_\ka=\la-w_0\la$ at $47/47$ injectives across
$A_3,A_4,A_5,D_4,D_5,D_6,D_7,E_6,E_7$;
$\dimv\soc_\sigma(I_\ka)=v$ for $706/706$ elements; the minuscule statement for
$256/256$; Proposition \ref{prop:lemA} for $981/981$ elements, of which $681$ are of type
$E$ and the rest of type $D$.
\item \texttt{socdim.py} --- Theorem \ref{thm:main} on $1149$ elements and
$12829$ individual socle steps, no exceptions.
\item \texttt{sweepall.py} --- the aggregate criterion-versus-solver sweep
quoted above; \texttt{DKKCAP} raises the orbit-size cutoff and \texttt{DKKP}
changes the characteristic.
\item \texttt{verify.py} --- the printed posets of \cite[pp.~82--83]{DKK}
reproduced with their Hasse edges; the criterion against an independent
graded-submodule solver; word-independence; Theorem \ref{thm:typeD}.
\item \texttt{trunctest.py} --- Proposition \ref{prop:T} (top truncation) on
$1000$ elements of $D_4$--$D_7$, every non-spin fundamental weight
($k\le n-2$) with orbit at most $400$, no exceptions.
\item \texttt{inwardlemma.py} --- Lemma \ref{lem:I}, exhaustively in $E_6$,
$E_7$, $E_8$ ($315$ instances) and for $D_4$--$D_{10}$ ($126$ instances).
\item \texttt{inwardproof.py} --- the inputs to Lemma \ref{lem:Id}: simple
socle in the top degree ($39$ injectives, $D_4$--$D_9$), multiplicity-freeness
on the chain for $\ka=1$, and the relation at a chain vertex ($550$ instances).
\item \texttt{bigrade.py} --- Lemma \ref{lem:bigrade}: $\Gamma$ is a tree and
the potential $f$ exists, for $A_2$--$A_8$, $D_4$--$D_9$, $E_6$, $E_7$, $E_8$.
\item \texttt{master.py} --- Proposition \ref{conj:staircase} over all $1124$
graded submodules of $I_\ka$ for $D_4$--$D_7$: exactly $78$ non-staircases,
every one with $\delta=2$.
\item \texttt{starchar.py} --- $(*)^{\circ}$ at $1998/1998$ extremal modules
over $D_5$--$D_{10}$, with all $111$ up-set failures at $t-1$ having $c_t=0$.
\item \texttt{localstar.py} --- the failure of the localisation of
$(*)^{\circ}$: $947/951$ of the graded submodules of $I_\ka$ satisfying its
hypotheses over $D_4$--$D_7$, the four exceptions all at $\ka=4$.
\item \texttt{reduce.py} --- Proposition \ref{prop:reps} and $(B)$: $15364$
instances over $D_4$--$D_{10}$, $k\le6$, collapsing to $83$ representatives;
reduction step $1705/1705$, $(B)$ at $43/43$.
\item \texttt{canon.py} --- $\tau$, the canonical truncation $T(v)$,
(T0)--(T2), and the classification of the closure obligations.
\item \texttt{strands.py} --- the strand decomposition (W1)--(W4), the closed
form, and the arithmetic steps of Lemma \ref{lem:death} and Proposition
\ref{prop:W4l}, over $D_5$--$D_{17}$ and every $\ka$.
\item \texttt{audit.py} --- the whole chain re-checked link by link, with
negative controls, characteristic- and word-independence, and a check that the
counts quoted here are the ones it computes.
\item \texttt{thresh.py} --- the closed forms for the profile of $I_\ka$ and
for $v_a$, checked against the modules, and the obligation classification over
$D_5$--$D_{60}$: $34104$ representatives, $26.2$ million obligations and
$2.99$ million instances of (T0)--(T2), all resolved.
\item \texttt{filltop.py}, \texttt{claimG.py}, \texttt{epscrit.py},
\texttt{violators.py} --- the unconditional reductions (a), (b), (c) of Remark
\ref{rem:Tk3} and the structure of the non-staircases.
\item \texttt{squeeze.py}, \texttt{homological.py}, \texttt{headrule.py},
\texttt{lemA.py}, \texttt{eproj.py} --- the individual ingredients.
\end{itemize}

Aggregate agreement of Corollary \ref{cor:crit} with the independent solver is
$2179/2179$, broken down as type $A$ $214$, type $D$ $1232$, type $E$ $733$;
this is the output of \texttt{sweepall.py}, which enumerates every fundamental
weight of $A_3$--$A_6$, $D_4$--$D_7$, $E_6$, $E_7\om_1$ and $E_7\om_7$ whose
Weyl orbit has at most $400$ elements and compares the criterion against the
solver element by element.  All computations are over a finite prime field; the
whole suite was rerun at $p=5,11,13$ with identical results.

\begin{thebibliography}{99}

\bibitem{BK} P.~Baumann and J.~Kamnitzer, \emph{Preprojective algebras and MV
polytopes}, Represent.\ Theory \textbf{16} (2012), 152--188.

\bibitem{BKT} P.~Baumann, J.~Kamnitzer and P.~Tingley, \emph{Affine
Mirkovi\'c--Vilonen polytopes}, Publ.\ Math.\ IH\'ES \textbf{120} (2014),
113--205.

\bibitem{CB} W.~Crawley-Boevey, \emph{On the exceptional fibres of Kleinian
singularities}, Amer.\ J.\ Math.\ \textbf{122} (2000), 1027--1037.

\bibitem{DKK} A.~Dinkins, D.~Karpov and V.~Krylov, \emph{The quantum Hikita
conjecture via quasimaps}, arXiv:2608.16746.

\bibitem{DFKKMSSVW} A.~Dranowski et al., \emph{Heaps, crystals, and
preprojective algebra modules}, arXiv:2202.02490.

\bibitem{GLS} C.~Geiss, B.~Leclerc and J.~Schr\"oer, \emph{Kac-Moody groups
and cluster algebras}, Adv.\ Math.\ \textbf{228} (2011), 329--433.

\bibitem{Nak} H.~Nakajima, \emph{Quiver varieties and Kac-Moody algebras},
Duke Math.\ J.\ \textbf{91} (1998), 515--560.

\bibitem{ST} A.~Savage and P.~Tingley, \emph{Quiver Grassmannians, quiver
varieties and the preprojective algebra}, Pacific J.\ Math.\ \textbf{251}
(2011), 393--429.

\bibitem{Stem} J.~Stembridge, \emph{Minuscule elements of Weyl groups}, J.\
Algebra \textbf{235} (2001), 722--743.

\end{thebibliography}

\end{document}
